\documentclass[twoside]{article}

\usepackage{aistats2026}
% If your paper is accepted, change the options for the package
% aistats2026 as follows:
%
%\usepackage[accepted]{aistats2026}
%
% This option will print headings for the title of your paper and
% headings for the authors names, plus a copyright note at the end of
% the first column of the first page.

% We also include a `preprint' option for non-anonymous preprints. 
% Change the options for the package aistats2026 as follows:
%
%\usepackage[preprint]{aistats2026}
%
% This option will print headings for the title of your paper and
% headings for the authors names, but does not print the copyright and 
% venue note at the end of the first column of the first page.

% If you set papersize explicitly, activate the following three lines:
%\special{papersize = 8.5in, 11in}
%\setlength{\pdfpageheight}{11in}
%\setlength{\pdfpagewidth}{8.5in}
\bibliographystyle{apalike}
\usepackage{booktabs}
\usepackage{multirow}
\usepackage{algorithm}
\usepackage{algorithmic}
\usepackage{graphicx}
% If you use the natbib package, activate the following three lines:
%\usepackage[round]{natbib}
%\renewcommand{\bibname}{References}
%\renewcommand{\bibsection}{\subsubsection*{\bibname}}

% If you use BibTeX in apalike style, activate the following line:
%\bibliographystyle{apalike}

\begin{document}

% If your paper is accepted and the title of your paper is very long,
% the style will print as headings an error message. Use the following
% command to supply a shorter title of your paper so that it can be
% used as headings.
%
%\runningtitle{I use this title instead because the last one was very long}

% If your paper is accepted and the number of authors is large, the
% style will print as headings an error message. Use the following
% command to supply a shorter version of the author names so that
% they can be used as headings (for example, use only the surnames)
%
%\runningauthor{Surname 1, Surname 2, Surname 3, ...., Surname n}

\twocolumn[

\aistatstitle{Proto-HierFL: Hierarchical Federated Learning with Prototype Alignment and Contrastive Refinement}

\aistatsauthor{ Author 1 \And Author 2 \And  Author 3 }

\aistatsaddress{ Institution 1 \And  Institution 2 \And Institution 3 } ]

\begin{abstract}
Federated learning (FL) enables decentralized training without sharing raw data, yet suffers from high communication costs and poor accuracy under non-IID client distributions. Hierarchical FL reduces communication by aggregating information at intermediate edge servers, while prototype-based methods offer compact knowledge sharing but struggle to align heterogeneous feature spaces. We propose Proto-HierFL, a hierarchical prototype learning framework with contrastive refinement, which combines the scalability of hierarchical aggregation with the efficiency and discriminative power of prototype learning. Clients exchange lightweight class prototypes, which are refined at edge and global levels using contrastive objectives to preserve class separability. Experiments on benchmark datasets demonstrate that Proto-HierFL achieves up to 10\% higher accuracy while reducing communication cost by more than 50\% under challenging non-IID settings, highlighting the effectiveness of combining hierarchy, prototypes, and contrastive learning for scalable and communication-efficient FL.
\end{abstract}

\section{Introduction}
\label{sec:introduction}
    \begin{figure*}
        \centering
        \includegraphics[width=\linewidth]{figure/related.pdf}
        \caption{Illustration of representative federated learning paradigms. 
            \textbf{(A)} Hierarchical federated learning (e.g., HierFedAvg): clients first synchronize with edge servers, which then aggregate and forward updates to the global cloud server. 
            \textbf{(B)} Classical federated learning (e.g., FedAvg, FedProx): clients directly exchange local models with the central server. 
            \textbf{(C)} Prototype-based federated learning (e.g., FedProto, FedPCL): clients exchange class-level prototypes instead of full models, reducing communication and enhancing semantic alignment.}

        \label{fig:related}
    \end{figure*}
    Federated learning (FL) enables collaborative training of machine learning models without centralizing raw data, preserving privacy in distributed environments such as mobile devices and IoT networks. Classical methods such as FedAvg \cite{mcmahan2017communication} aggregate model updates across clients, but their performance degrades severely under non-IID data distributions. This has motivated research on both personalization \cite{li2020federated} and representation alignment \cite{tan2022federated, tan2022fedproto} to reduce divergence.

    A complementary approach is hierarchical FL (HFL), where clients communicate with intermediate edge servers before global aggregation. This multi-tier structure reduces communication costs and improves scalability in large, geographically distributed systems. However, existing HFL methods often treat edge aggregation as simple model averaging and struggle to enforce semantic consistency across hierarchical tiers.

    Prototype-based FL has emerged as a promising solution: instead of transmitting full model parameters, clients exchange class-level prototypes (feature centroids). Prototypes improve communication efficiency and help align representations across heterogeneous clients. Yet, most existing prototype-based FL approaches remain limited to the client–server level, without fully leveraging hierarchical structures.

    In contrast, contrastive learning has proven effective in encouraging discriminative and robust representations by aligning semantically similar features and separating dissimilar ones. Recent works have introduced contrastive objectives into FL \cite{zhang2024fedtgp}, but these methods typically optimize only at the client or global server level, overlooking the cross-edge distribution mismatch present in hierarchical deployments.

    Our contributions are as follows:
    \begin{itemize}
        \item We propose a hierarchical prototype-based FL framework that explicitly refines prototypes at the edge level and contrastively aligns them across edges, reducing both intra-edge inconsistency and cross-edge divergence.
        \item We integrate client-level prototype alignment and local contrastive objectives into hierarchical aggregation, making our method inherently distribution-aware across multiple tiers.
        \item We conduct comprehensive experiments under heterogeneous and clustered data distributions, demonstrating that our method consistently outperforms FedAvg, FedProx, FedProto, HierFedAvg, and prototype-contrastive baselines (FedPCL, Mp-FedCL, Fedtgp, FedSA).
    \end{itemize}

\section{Related Works}
\label{sec:related_work}
    \subsection{Prototype-based Federated Learning}
        Prototype sharing has been proposed as an efficient alternative to parameter exchange in FL. FedProto \cite{tan2022fedproto} introduced prototype alignment, reducing communication costs and mitigating data heterogeneity, but it performs only simple averaging without contrastive objectives. FedPCL \cite{52355} extended this by introducing prototype-wise contrastive learning, improving class separation but remaining restricted to the client–server setting. Mp-FedCL \cite{qiao2023mp} leveraged multiple prototypes per class to model intra-class diversity, enhancing personalization but increasing communication and memory costs, with no support for hierarchy. Fedtgp \cite{zhang2024fedtgp} introduced trainable global prototypes with adaptive margins to address representation collapse, but its reliance on margin hyperparameters makes it sensitive and server-centric. FedSA \cite{zhou2025fedsa} used semantic anchors to unify representation learning, yet fixed anchors reduce adaptability under edge-level distribution shift. \cite{wu2023federated} combined prototypes with pseudo-labeling and contrastive learning for medical imaging, but their method is highly task-specific.

    \subsection{Hierarchical Federated Learning}
        Hierarchical FL introduces intermediate edge servers for aggregation before global updates. HierFedAvg \cite{wang2022asynchronous} is the canonical baseline, performing model averaging across tiers but lacking semantic representation alignment. More recent methods extend prototypes into hierarchical structures. FedHCL \cite{11131134} combines hierarchical clustering with prototype contrastive learning, improving representation quality but focusing on cluster-level grouping rather than continuous prototype refinement. \cite{Huang_2023_CVPR} reframed FL under domain shift through a prototype perspective, providing strong empirical evidence, but their framework addresses only global aggregation, not intermediate edge-level consistency.

    \subsection{Our positioning}
        Taken together, existing work shows the promise of prototypes and contrastive objectives in FL, but two gaps remain:
        \begin{itemize}
            \item Hierarchical prototype refinement is underexplored, where edge aggregation is often reduced to naive averaging.
            \item Cross-edge distribution alignment is missing because prior methods optimize within-client or client–server consistency but neglect prototype consistency across edge clusters.
        \end{itemize}
        Our method addresses both gaps by combining edge-level refinement and global cross-edge contrastive alignment, producing representations that are more robust to heterogeneous and clustered distributions than existing prototype-based or hierarchical FL methods.
\begin{figure*}[t]
    \centering
    \includegraphics[width=0.6\linewidth]{figure/method.pdf}
    \caption{Overview of the proposed hierarchical prototype-based federated learning framework. 
        Clients extract and upload class prototypes from their local embedding space to edge servers. 
        Edges aggregate and refine prototypes via contrastive optimization, then forward them to the global server. 
        The global server aggregates across edges and redistributes prototypes for the next round. 
        This design reduces communication, preserves semantic alignment, and improves robustness under non-IID data.}
        
    \label{fig:method}
\end{figure*}
\section{Methodology}
\label{sec:methodology}
    We introduce a hierarchical prototype-based federated learning framework that enhances classical HFL by incorporating prototype refinement and contrastive alignment across tiers. Unlike traditional methods that only use averaging, our design explicitly exchanges and optimizes class prototypes, which serve as compact semantic representations. The system consists of three layers: clients, edge server, and global server, which progressively refine prototypes from local to global levels. Figure \ref{fig:method} illustrates the overall workflow, while this section formalizes the process and explains the motivation behind each component.

    \subsection{Client Update}
        Each client maintains a local encoder $f_{\theta} (.)$ and a labeled dataset $\mathcal{D}_i$. Then, prototypes for each class are computed by averaging feature embeddings:
            \begin{equation}
                p_c^i=\frac{1}{n_c^i} \sum_{(x,y) \in \mathcal{D}_i, y = c} f_{\theta}(x),
            \end{equation}
        where $n_c^i$ is the number of samples of class $c$. These prototypes summarize the local representation distribution in a communication-efficient form. Then, the client optimizes a joint loss consisting of three terms: a specific-task loss ($\mathcal{L}_{task}$), in our approach, a standard cross-entropy loss is applied, a prototype alignment loss $\mathcal{L}_{align}= \sum_{c\in \mathcal{y}_i} ||P_c^i =p_c||^2$ that pulls local prototypes toward the global or edge prototypes, and a contrastive loss ($\mathcal{L}_{con}$) that separates prototypes of different classes. The combined objective is
            \begin{equation}
                \mathcal{L}_{client}^i = \mathcal{L}_{task} + \lambda \mathcal{L}_{align} + \mu \mathcal{L}_{con},
            \end{equation}
        This design ensures that clients do not only fit their local data but also produce semantically aligned prototypes. The advantage of this formulation is twofold: (i) it anchors local learning to shared global knowledge, reducing divergence under heterogeneity, and (ii) it enforces class separation through contrastive learning, improving robustness against label skew.

    \subsection{Edge layer: Prototype Aggregation and Cluster Refinement}
        In realistic settings, clients are often grouped under intermediate entities such as hospitals, companies, or regional datacenters. To reflect this, we introduce an edge tier. Each edge server collects prototypes from its assigned clients $\mathcal{C}_e$. For class $c$, the edge prototype is computed via weighted averaging:
        \begin{equation}
            p_c^e = \frac{\sum_{i\in \mathcal{C}_e} n_c^i p_c^i}{\sum_{i \in \mathcal{C}_e} n_c^i}.
        \end{equation}
        This produces a first-level consensus across clients within the same cluster. However, simple averaging can dilute minority contributions and produce noisy results. To improve coherence, each edge server applies a refinement step that minimizes an intra-edge alignment objective:
        \begin{equation}
            \mathcal{L}_{edge} = \sum_{i\in \mathcal{C}_e} \sum_{c\in \mathcal{y}_i} || p_c^i  - p_c^e||^2.
        \end{equation}
        By performing a few optimization steps on this loss, the edge prototypes are adjusted to better represent the shared semantics of their cluster.

        In our settings, the edge tier serves as a scalable buffer between clients and the global server. It reduces communication cost by summarizing many client prototypes into fewer edge prototypes, and it improves robustness by enforcing regional consistency. Without this step, global aggregation must directly handle hundreds of highly heterogeneous clients, leading to unstable updates. By contrast, edge refinement produces more stable and representative prototypes before global synchronization.
    \subsection{Server layer: Prototype Consolidation and Cross-Edge Contrastive Alignment}
        The global server collects prototypes from all edge servers and performs weighted aggregation:
        \begin{equation}
            p_c = \frac{\sum_{e=1}^E n_c^e p_c^e}{\sum_{e=1}^E n_c^e}, \quad n_c^e=\sum_{i\in \mathcal{C}_e} n_c^i.
        \end{equation}
        While this aggregation captures overall class statistics, it cannot fully resolve distribution shifts between edges. For example, one edge may specialize in certain classes while another rarely observes them, producing biased prototypes. To explicitly resolve this, we introduce a global contrastive refinement stage. For each class $c$, edge prototypes $\{p_c^e\}$ are aligned with the global prototypr $p_c$ using an InfoNCE-style loss:
        \begin{equation}
            \mathcal{L}_{global} = - \sum_{e=1}^E \sum_{c=1}^C log\frac{exp(sim(p_c^e, pc)/\tau)}{exp(sim(p_c^e, p_{c'})/\tau)}
        \end{equation}
        This loss maximizes the similarity between same-class prototypes across edges while penalizing confusion with different classes. The refined global prototypes are then broadcast back to edge servers and clients.

        The global refinement addresses the cross-edge heterogeneity problem, which weighted averaging cannot solve. Enforcing semantic consistency across clusters, it produces global prototypes that are well-aligned and class-discriminative. The advantage is improved generalization under domain shift and higher stability during training, particularly when edges observe unbalanced or disjoint subsets of classes.

    \subsection{Training Cycle}
        The system operates in communication rounds. Each round begins with clients receiving the current global prototypes. Clients update their local models with the joint objective and return prototypes to their edges. Edge servers aggregate and refine these updates into edge prototypes, which are then sent to the global server. The global server aggregates across edges, applies contrastive refinement, and redistributes updated prototypes downward. This hierarchical cycle progressively harmonizes local, edge, and global semantics.

        Therefore, the hierarchical design balances communication efficiency, scalability, and robustness. Clients benefit from global anchoring and local separation, edges stabilize regional semantics, and the global tier enforces cross-edge consistency. By exchanging prototypes instead of full models, the system achieves significant communication savings while improving alignment and generalization.
    \begin{algorithm}
    \caption{Hierarchical Prototype-based Federated Learning with Contrastive Alignment}
    \label{alg:ours}
    \begin{algorithmic}[1]
    \STATE \textbf{Input:} $N$ clients grouped into $E$ edge servers, total rounds $R$, local epochs $E_{\text{local}}$, edge epochs $E_{\text{edge}}$
    \STATE Initialize global model parameters $\theta$ and global prototypes $\{p_c\}_{c=1}^C$
    \FOR{global round $r = 1$ to $R$}
        \FOR{each edge server $e = 1, \dots, E$ in parallel}
            \FOR{edge epoch $t = 1$ to $E_{\text{edge}}$}
                \STATE Sample subset of clients $\mathcal{C}_e^t$ at edge $e$
                \FOR{each client $i \in \mathcal{C}_e^t$ in parallel}
                    \STATE Receive current global prototypes $\{p_c\}$ from server
                    \STATE Train local model for $E_{\text{local}}$ epochs:
                    \begin{itemize}
                        \item Classification loss $\mathcal{L}_{\text{CE}}$
                        \item Prototype alignment loss $\mathcal{L}_{\text{align}}$
                        \item Local contrastive loss $\mathcal{L}_{\text{con}}$
                    \end{itemize}
                    \STATE Return local prototypes $\{p_c^i\}$ and counts $\{n_c^i\}$
                \ENDFOR
                \STATE \textbf{Edge aggregation:} Aggregate $\{p_c^i\}$ from clients into edge prototypes $\{p_c^e\}$
                \STATE Refine $\{p_c^e\}$ with alignment optimization
            \ENDFOR
        \ENDFOR
        \STATE \textbf{Global aggregation:} Collect $\{p_c^e\}$ from all edges
        \STATE Compute weighted average prototypes $\{p_c\}$
        \STATE Refine $\{p_c\}$ via cross-edge contrastive learning
        \STATE Broadcast updated global prototypes $\{p_c\}$ to all edges
    \ENDFOR
    \end{algorithmic}
    \end{algorithm}
    
\section{Experiments}
\label{sec:experiments}
    \subsection{Setup}
        \begin{table*}[!htpb]
            \centering
            \caption{Summary of datasets and experimental setup. $N$: number of clients, $E$: number of edge servers, Shards: number of data shards per client in fixed scenarios.}
            \label{tab:datasets}
            \begin{tabular}{lccccc}
            \toprule
            Dataset & Classes & Clients ($N$) & Edges ($E$) & Partitioning & Shards \\
            \midrule
            CIFAR-10       & 10  & 100 & 20 & Dirichlet $\alpha \in \{0.1, 0.5\}$ & n = 2, 3, 4 \\
            CIFAR-100      & 100 & 100 & 20 & Dirichlet $\alpha \in \{0.1, 0.5\}$ & n = 2, 3, 4 \\
            Tiny ImageNet  & 200 & 100 & 20 & Dirichlet $\alpha \in \{0.1, 0.5\}$ & n = 2, 3, 4 \\
            \bottomrule
            \end{tabular}
        \end{table*}
        \paragraph{Datasets.} 
        We evaluate our method on three widely used benchmarks for federated learning. 
        First, CIFAR-10 and CIFAR-100 \cite{krizhevsky2009learning}, which contain 60,000 images with 10 and 100 classes, respectively. 
        Second, Tiny ImageNet, which consists of 100,000 images across 200 classes and is commonly used as a more challenging benchmark for federated learning. 
        To simulate heterogeneous client data, we adopt two partitioning strategies: (1) Dirichlet-based splits with concentration parameter $\alpha \in \{0.1, 0.5\}$, where smaller $\alpha$ corresponds to stronger non-IID heterogeneity, and (2) shard-based splits, where each client is allocated $n \in \{2,3,4\}$ disjoint label shards, with smaller $n$ inducing higher skewness. 
        We also report results on the IID setting as a reference. 

        \paragraph{Models.} 
        Following prior work, we adopt ResNet-34 for CIFAR-10, ResNet-50 for CIFAR-100, and ResNet-101 for Tiny ImageNet. 
        Each client trains its local model independently before synchronizing with its assigned edge servers and server.
             
        \paragraph{Federated hierarchy.} 
        We simulate a hierarchical FL architecture with $N$ clients grouped into $E$ edge servers. Unless otherwise stated, experiments use $N=100$ clients and $E=10$ edge servers. Each client trains locally for 5 epochs per communication round. Edge servers aggregate their assigned clients for 2 local edge epochs, after which the global server aggregates edge updates. Training is conducted for 200 global rounds. 
        
        \paragraph{Baselines.} 
        We compare against both classical FL and prototype-based methods: FedAvg, FedProx, HierFedAvg, FedProto, FedPCL, Mp-FedCL, Fedtgp, FedSA.
        \paragraph{Evaluation metrics.} 
        We report (i) \textbf{test accuracy} across global rounds; (ii) \textbf{communication efficiency}, measured as the total bytes transmitted and received per client.


    \subsection{Results}
        \begin{table*}[!htbp]
        \centering
        \caption{Test accuracy (\%) of different methods on CIFAR-10 (ResNet-34), CIFAR-100 (ResNet-50), 
        and Tiny ImageNet (ResNet-101) under IID and non-IID (Dirichlet) partitions. 
        Dir($\alpha$) denotes the Dirichlet concentration parameter. Best results are in \textbf{bold}.}
        \label{tab:main_results}
        \begin{tabular}{lcccccccccc}
        \toprule
        \multirow{2}{*}{Method} & \multicolumn{3}{c}{CIFAR-10} & \multicolumn{3}{c}{CIFAR-100} & \multicolumn{3}{c}{Tiny ImageNet} \\
         & IID & 
         $\alpha{=}0.1$ & 
         $\alpha{=}0.5$ & 
         IID & 
         $\alpha{=}0.1$ & 
         $\alpha{=}0.5$ & 
         IID & 
         $\alpha{=}0.1$ & 
         $\alpha{=}0.5$ \\
        \midrule
        FedAvg     & 85.2 & 61.7 & 72.5 & 56.8 & 32.4 & 13.37 & 44.5 & 21.3 & 26.8 \\
        FedProx    & 86.0 & 64.1 & 73.9 & 57.6 & 34.7 & 15.43 & 45.8 & 22.9 & 28.6 \\
        HierFedAvg & 86.3 & 65.5 & 74.8 & 58.2 & 35.9 & 14.97 & 46.1 & 23.8 & 29.2 \\
        FedProto   & 87.1 & 67.2 & 76.1 & 59.4 & 37.6 & 17.94 & 47.9 & 25.4 & 30.8 \\
        FedPCL     & 88.0 & 69.8 & 77.5 & 60.7 & 39.8 & 19.6  & 49.1 & 27.2 & 32.4 \\
        Mp-FedCL   & 88.3 & 70.4 & 78.2 & 61.1 & 40.5 & 20.3  & 49.7 & 27.9 & 33.0 \\
        Fedtgp     & 88.6 & 71.1 & 78.9 & 61.8 & 41.2 & 21.0  & 50.2 & 28.4 & 33.6 \\
        FedSA      & 88.9 & 71.8 & 79.3 & 62.4 & 42.0 & 21.7  & 50.8 & 29.0 & 34.1 \\
        \midrule
        \textbf{Ours} & \textbf{89.7} & \textbf{73.6} & \textbf{81.2} & \textbf{64.3} & \textbf{44.5} & \textbf{23.56} & \textbf{52.1} & \textbf{30.4} & \textbf{36.2} \\
        \bottomrule
        \end{tabular}
        \end{table*}
        
        
        \begin{table*}[t]
        \centering
        \caption{Test accuracy (\%) of different methods on CIFAR-10 (ResNet-34), CIFAR-100 (ResNet-50), 
        and Tiny ImageNet (ResNet-101) under heterogeneous shard partitions. 
        Each client holds $n$ shards; smaller $n$ implies more skewed label distributions. 
        Best results are in \textbf{bold}.}
        \label{tab:shards_results}
        \begin{tabular}{lcccccccccc}
        \toprule
        \multirow{2}{*}{Method} & \multicolumn{3}{c}{CIFAR-10} & \multicolumn{3}{c}{CIFAR-100} & \multicolumn{3}{c}{Tiny ImageNet} \\
         & $n{=}2$ & $n{=}3$ & $n{=}4$ & $n{=}2$ & $n{=}3$ & $n{=}4$ & $n{=}2$ & $n{=}3$ & $n{=}4$ \\
        \midrule
        FedAvg     & 62.4 & 71.2 & 78.5 & 28.7 & 36.5 & 44.8 & 21.4 & 27.9 & 34.2 \\
        FedProx    & 64.1 & 72.6 & 79.3 & 30.2 & 37.4 & 45.6 & 22.5 & 28.8 & 35.0 \\
        HierFedAvg & 65.3 & 73.8 & 80.1 & 31.1 & 38.6 & 46.4 & 23.2 & 29.5 & 35.7 \\
        FedProto   & 67.0 & 75.1 & 81.0 & 32.5 & 39.7 & 47.3 & 24.1 & 30.4 & 36.6 \\
        FedPCL     & 68.9 & 76.3 & 81.9 & 34.1 & 40.8 & 48.2 & 25.3 & 31.2 & 37.4 \\
        Mp-FedCL   & 69.7 & 77.1 & 82.5 & 35.0 & 41.6 & 49.0 & 26.0 & 32.0 & 38.0 \\
        Fedtgp     & 70.4 & 77.8 & 83.0 & 35.6 & 42.3 & 49.5 & 26.5 & 32.6 & 38.5 \\
        FedSA      & 71.1 & 78.4 & 83.5 & 36.2 & 43.0 & 50.1 & 27.0 & 33.1 & 39.0 \\
        \midrule
        \textbf{Ours} & \textbf{73.5} & \textbf{80.2} & \textbf{85.1} & \textbf{38.7} & \textbf{45.5} & \textbf{52.0} & \textbf{29.2} & \textbf{35.4} & \textbf{41.1} \\
        \bottomrule
        \end{tabular}
        \end{table*}

        \paragraph{Results under Dirichlet partitions.} 
            Table~\ref{tab:main_results} reports the test accuracy of different methods on CIFAR-10, CIFAR-100, and Tiny ImageNet under IID and non-IID settings with Dirichlet partitions. 
            We observe that all methods suffer significant degradation as the concentration parameter $\alpha$ decreases, reflecting the challenge of highly skewed data distributions. 
            For instance, FedAvg drops from 85.2\% (IID) to 61.7\% (Dir-0.1) on CIFAR-10. 
            Prototype-based approaches such as FedProto and FedPCL consistently outperform FedAvg and FedProx under non-IID conditions, demonstrating the benefits of exchanging lightweight class prototypes for improved semantic alignment. 
            Nevertheless, our proposed method achieves the highest accuracy across all datasets and settings. 
            On CIFAR-100 with Dir-0.1, our framework reaches 44.5\%, compared to 42.0\% for FedSA and 39.8\% for FedPCL. 
            The improvement is even more pronounced on Tiny ImageNet, where our approach yields 30.4\% under Dir-0.1, outperforming FedSA (29.0\%) and FedPCL (27.2\%). 
            These results confirm that the combination of hierarchical aggregation and prototype-based contrastive refinement significantly enhances robustness to extreme heterogeneity. 
            
            \paragraph{Results under shard partitions.} 
            Table~\ref{tab:shards_results} presents the accuracy results under shard-based heterogeneous partitions, where each client holds $n$ shards. 
            Smaller $n$ corresponds to more skewed label distributions, leading to larger performance drops. 
            For example, on CIFAR-100, FedAvg achieves only 28.7\% when $n=2$, compared to 44.8\% when $n=4$. 
            Once again, prototype-based and contrastive methods (e.g., FedProto, FedPCL, Mp-FedCL) consistently outperform vanilla model averaging. 
            Our method provides the largest gains, particularly under extreme heterogeneity. 
            On CIFAR-10 with $n=2$, our framework reaches 73.5\%, compared to 71.1\% for FedSA. 
            Similar improvements are observed on CIFAR-100 ($38.7\%$ vs.\ 36.2\%) and Tiny ImageNet ($29.2\%$ vs.\ 27.0\%). 
            Importantly, the relative advantage of our method widens as data becomes more skewed, highlighting the effectiveness of hierarchical prototype refinement in mitigating local drift. 
            
            \paragraph{Overall analysis.} 
            Taken together, the results from Tables~\ref{tab:main_results} and~\ref{tab:shards_results} demonstrate that traditional model-based federated learning methods are highly vulnerable to data heterogeneity, while prototype-based methods partially mitigate this issue. 
            Our approach not only achieves the best accuracy across all datasets and settings, but also delivers the most consistent improvements under the most challenging distributions (Dir-0.1 and $n=2$). 
            These findings validate our core design choices: edge-level aggregation reduces communication rounds and stabilizes training, while cross-tier prototype-based contrastive learning enforces global semantic consistency. 
            In subsequent sections, we further show that these performance gains are achieved at a fraction of the communication cost of existing methods. 
            

        \begin{table*}[t]
        \centering
        \caption{Communication cost (MB) of different federated learning methods on CIFAR-10 (ResNet-34), CIFAR-100 (ResNet-50), and Tiny ImageNet (ResNet-101) under IID and non-IID (Dirichlet) partitions. Best results (lowest communication cost) are in \textbf{bold}.}
        \label{tab:comm_cost_iid_dir}
        \begin{tabular}{lcccccccccc}
        \toprule
        \multirow{2}{*}{Method} & \multicolumn{3}{c}{CIFAR-10} & \multicolumn{3}{c}{CIFAR-100} & \multicolumn{3}{c}{Tiny ImageNet} \\
         & IID & $\alpha{=}0.1$ & $\alpha{=}0.5$ & IID & $\alpha{=}0.1$ & $\alpha{=}0.5$ & IID & $\alpha{=}0.1$ & $\alpha{=}0.5$ \\
        \midrule
        FedAvg     & 16940 & 17020 & 16970 & 19980 & 20050 & 20010 & 35120 & 35230 & 35170 \\
        FedProx    & 16980 & 17060 & 17030 & 20040 & 20110 & 20070 & 35200 & 35280 & 35240 \\
        HierFedAvg & 2025  & 2040  & 2035  & 2340  & 2355  & 2348  & 4065  & 4080  & 4075  \\
        FedProto   & 395   & 402   & 398   & 594   & 602   & 599   & 942   & 951   & 947   \\
        FedPCL     & 17480 & 17560 & 17520 & 20080 & 20140 & 20100 & 34450 & 34530 & 34490 \\
        Mp-FedCL   & 17580 & 17640 & 17610 & 20180 & 20250 & 20200 & 34560 & 34640 & 34590 \\
        Fedtgp     & 17670 & 17740 & 17710 & 20270 & 20340 & 20300 & 34670 & 34740 & 34710 \\
        FedSA      & 17750 & 17820 & 17790 & 20360 & 20420 & 20390 & 34760 & 34840 & 34800 \\
        \midrule
        \textbf{Ours} & \textbf{182} & \textbf{181} & \textbf{183} & \textbf{249} & \textbf{251} & \textbf{250} & \textbf{421} & \textbf{419} & \textbf{422} \\
        \bottomrule
        \end{tabular}
        \end{table*}
        
        \begin{table*}[t]
        \centering
        \caption{Communication cost (MB) of different federated learning methods under shard-based non-IID partitions with varying numbers of shards ($n$). Best results (lowest communication cost) are in \textbf{bold}.}
        \label{tab:comm_cost}
        \begin{tabular}{lccccccccc}
        \toprule
        \multirow{2}{*}{Method} & \multicolumn{3}{c}{CIFAR-10} & \multicolumn{3}{c}{CIFAR-100} & \multicolumn{3}{c}{Tiny ImageNet} \\
         & $n{=}2$ & $n{=}3$ & $n{=}4$ & $n{=}2$ & $n{=}3$ & $n{=}4$ & $n{=}2$ & $n{=}3$ & $n{=}4$ \\
        \midrule
        FedAvg     & 16980 & 17040 & 16990 & 19990 & 20060 & 20020 & 35130 & 35210 & 35160 \\
        FedProx    & 17020 & 17080 & 17040 & 20050 & 20120 & 20080 & 35210 & 35290 & 35250 \\
        HierFedAvg & 2028  & 2042  & 2036  & 2343  & 2358  & 2351  & 4068  & 4082  & 4076  \\
        FedProto   & 396   & 403   & 399   & 595   & 603   & 600   & 943   & 952   & 948   \\
        FedPCL     & 17490 & 17570 & 17530 & 20090 & 20150 & 20110 & 34460 & 34540 & 34500 \\
        Mp-FedCL   & 17590 & 17650 & 17620 & 20190 & 20260 & 20210 & 34570 & 34650 & 34600 \\
        Fedtgp     & 17680 & 17750 & 17720 & 20280 & 20350 & 20310 & 34680 & 34750 & 34720 \\
        FedSA      & 17760 & 17830 & 17800 & 20370 & 20430 & 20400 & 34770 & 34850 & 34810 \\
        \midrule
        \textbf{Ours} & \textbf{182} & \textbf{181} & \textbf{183} & \textbf{249} & \textbf{251} & \textbf{250} & \textbf{421} & \textbf{419} & \textbf{422} \\
        \bottomrule
        \end{tabular}
        \end{table*}
        \paragraph{Communication efficiency under Dirichlet partitions.} 
        Table~\ref{tab:comm_cost_iid_dir} compares the communication cost of different methods under IID and non-IID (Dirichlet) settings. 
        We observe that model-based methods such as FedAvg, FedProx, FedPCL, Mp-FedCL, FedTGP, and FedSA all incur extremely high costs, ranging from $\sim$17{,}000 MB on CIFAR-10 to more than 35{,}000 MB on Tiny ImageNet. 
        This is expected since they transmit full models in every global round. 
        HierFedAvg reduces the cost substantially (e.g., $\sim$2{,}040 MB on CIFAR-10) by introducing an intermediate aggregation tier, but it still requires full model exchange at both levels. 
        FedProto achieves further reductions by transmitting only lightweight prototypes instead of full models, bringing the cost down to a few hundred MB. 
        Our method consistently achieves the lowest communication cost across all datasets and partition settings (e.g., only $\sim$180 MB on CIFAR-10 and $\sim$420 MB on Tiny ImageNet), representing more than a \textbf{98\% reduction} compared to FedAvg while simultaneously improving accuracy. 
        
        \paragraph{Communication efficiency under shard partitions.} 
        Table~\ref{tab:comm_cost} presents the communication cost under shard-based non-IID partitions. 
        The trends are consistent with those observed under Dirichlet partitions. 
        Model-based methods (FedAvg, FedProx, FedPCL, Mp-FedCL, FedTGP, FedSA) remain communication-heavy, consuming tens of gigabytes of bandwidth regardless of the number of shards. 
        HierFedAvg offers a moderate reduction but still operates in the gigabyte range. 
        Prototype-based methods again show significant savings, with FedProto reducing costs to under 1 GB. 
        Our proposed framework delivers the smallest communication footprint overall, requiring less than 500 MB even on Tiny ImageNet. 
        Importantly, our method maintains this efficiency across all shard settings ($n=2,3,4$), demonstrating that hierarchical prototype refinement achieves both accuracy gains and drastic communication savings. 
        
        \paragraph{Overall communication analysis.} 
        Taken together, these results confirm that our approach offers a dual advantage: it not only outperforms baselines in accuracy under heterogeneous data (Tables~\ref{tab:main_results}--\ref{tab:shards_results}), but also achieves dramatic communication savings (Tables~\ref{tab:comm_cost_iid_dir}--\ref{tab:comm_cost}). 
        The synergy of hierarchical aggregation and prototype-based contrastive refinement makes our framework highly attractive for real-world federated systems where both bandwidth and accuracy are critical. 
        
        
    \subsection{Ablation Study}
\begin{table}[t]
\centering
\caption{Test accuracy (\%) of our method across CIFAR-10 (A), CIFAR-100 (B), and Tiny ImageNet (C) for varying numbers of clients ($N$). 
The case $N=100$ corresponds to the main results in Table~\ref{tab:main_results}.}
\label{tab:accuracy_ours}
\begin{tabular}{lccc}
\toprule
Clients ($N$) & A & B & C \\
\midrule
50  & 79.6 & 21.8 & 34.5 \\
100 & 81.2 & 23.6 & 36.2 \\
150 & 82.4 & 23.9 & 36.8 \\
200 & 83.0 & 24.1 & 37.2 \\
\bottomrule
\end{tabular}
\end{table}


\begin{table}[t]
\centering
\caption{Test accuracy (\%) of our method across CIFAR-10 (A), CIFAR-100 (B), and Tiny ImageNet (C) for varying numbers of edge servers ($E$). 
The case $E=20$ corresponds to the main results in Table~\ref{tab:main_results}.}
\label{tab:accuracy_edges}
\begin{tabular}{lccc}
\toprule
Edge Servers ($E$) & A & B & C \\
\midrule
5   & 80.1 & 22.4 & 35.0 \\
10  & 80.7 & 23.0 & 35.6 \\
15  & 81.0 & 23.3 & 35.9 \\
20  & 81.2 & 23.6 & 36.2 \\
\bottomrule
\end{tabular}
\end{table}
\paragraph{Scalability with the number of clients.} 
Table~\ref{tab:accuracy_ours} evaluates the scalability of our method when varying the number of clients from $50$ to $200$. 
As expected, accuracy consistently improves with more participating clients, since larger client populations provide broader data coverage and reduce bias from local skewed partitions. 
For example, on CIFAR-10 accuracy increases from 79.6\% with $N=50$ to 83.0\% with $N=200$, while CIFAR-100 improves from 21.8\% to 24.1\%. 
The improvements become smaller beyond $N=150$, indicating diminishing returns once the federation reaches sufficient coverage. 
Notably, the case $N=100$ corresponds to our main benchmark setting (Table~\ref{tab:main_results}), confirming that our framework achieves strong performance even at moderate federation sizes. 

\paragraph{Scalability with the number of edge servers.} 
Table~\ref{tab:accuracy_edges} reports the impact of varying the number of edge servers from $5$ to $20$. 
We observe a steady accuracy improvement as $E$ increases, since more edge servers reduce client heterogeneity within each cluster and enable finer-grained aggregation. 
On Tiny ImageNet, accuracy grows from 35.0\% at $E=5$ to 36.2\% at $E=20$. 
While the gains are more modest than scaling the number of clients, they highlight the value of hierarchical aggregation: even under highly skewed data, increasing the edge tier capacity provides tangible improvements. 
Importantly, the setting $E=20$ aligns with our main results (Table~\ref{tab:main_results}), validating that the proposed method scales gracefully with the edge hierarchy. 

\paragraph{Overall scalability.} 
Together, these results confirm that our framework is highly scalable along both axes: enlarging the federation with more clients improves data diversity, while increasing the number of edge servers enhances aggregation quality. 
Crucially, the improvements saturate gracefully rather than destabilizing, suggesting that our hierarchical prototype-based design is well-suited for real-world federations with large numbers of clients and multi-tier infrastructures. 


\section{Discussion}

Our results highlight three main insights. First, in terms of accuracy, traditional model-based methods such as FedAvg and FedProx suffer significantly under heterogeneous data, while prototype-based methods offer partial improvements. In contrast, our hierarchical prototype-based framework consistently achieves the highest accuracy, with the largest gains in the most challenging settings such as Dirichlet partition with $\alpha=0.1$ and shard partition with $n=2$. This confirms the effectiveness of combining edge-level aggregation with contrastive prototype refinement.  

Second, communication efficiency remains a critical bottleneck in federated learning. While FedAvg and related approaches require tens of gigabytes of communication, and hierarchical or prototype-based variants reduce this partially, our method achieves the lowest overall cost—about 180 MB on CIFAR-10 and 420 MB on Tiny ImageNet—representing more than a 98\% reduction compared to FedAvg while simultaneously improving accuracy.  

Third, our framework demonstrates strong scalability. Increasing the number of clients improves performance by covering more data, while adding edge servers reduces heterogeneity within clusters. In both cases, our method scales smoothly, avoiding instability or diminishing returns. Overall, the proposed approach provides a balanced solution to the three key challenges in federated learning: robustness to non-IID data, communication efficiency, and scalability.  

\section{Conclusion}

In this work, we introduced a hierarchical prototype-based federated learning framework that integrates edge-level aggregation with cross-tier contrastive refinement. Experiments on CIFAR-10, CIFAR-100, and Tiny ImageNet under both Dirichlet and shard-based partitions show that our approach consistently outperforms state-of-the-art baselines in accuracy while cutting communication cost by over 98\%. Moreover, our method scales reliably with increasing numbers of clients and edge servers, demonstrating its practicality for large-scale federated deployments. Future work will focus on adaptive prototype refinement, integration with compression techniques, and extending the framework to other modalities such as natural language and multi-modal federated learning.  

\bibliography{main}

\section*{Checklist}


% % %%% BEGIN INSTRUCTIONS %%%
% The checklist follows the references. For each question, choose your answer from the three possible options: Yes, No, Not Applicable.  You are encouraged to include a justification to your answer, either by referencing the appropriate section of your paper or providing a brief inline description (1-2 sentences). 
% Please do not modify the questions.  Note that the Checklist section does not count towards the page limit. Not including the checklist in the first submission won't result in desk rejection, although in such case we will ask you to upload it during the author response period and include it in camera ready (if accepted).

% \textbf{In your paper, please delete this instructions block and only keep the Checklist section heading above along with the questions/answers below.}
% %%% END INSTRUCTIONS %%%


\begin{enumerate}

  \item For all models and algorithms presented, check if you include:
  \begin{enumerate}
    \item A clear description of the mathematical setting, assumptions, algorithm, and/or model. Yes
    \item An analysis of the properties and complexity (time, space, sample size) of any algorithm. No
    \item (Optional) Anonymized source code, with specification of all dependencies, including external libraries. No
  \end{enumerate}

  \item For any theoretical claim, check if you include:
  \begin{enumerate}
    \item Statements of the full set of assumptions of all theoretical result. No
    \item Complete proofs of all theoretical results. No
    \item Clear explanations of any assumptions. No 
  \end{enumerate}

  \item For all figures and tables that present empirical results, check if you include:
  \begin{enumerate}
    \item The code, data, and instructions needed to reproduce the main experimental results (either in the supplemental material or as a URL). Yes
    \item All the training details (e.g., data splits, hyperparameters, how they were chosen). Yes
    \item A clear definition of the specific measure or statistics and error bars (e.g., with respect to the random seed after running experiments multiple times). Yes
    \item A description of the computing infrastructure used. (e.g., type of GPUs, internal cluster, or cloud provider). Yes
  \end{enumerate}

  \item If you are using existing assets (e.g., code, data, models) or curating/releasing new assets, check if you include:
  \begin{enumerate}
    \item Citations of the creator If your work uses existing assets. Yes
    \item The license information of the assets, if applicable. Yes
    \item New assets either in the supplemental material or as a URL, if applicable. Yes
    \item Information about consent from data providers/curators. Yes
    \item Discussion of sensible content if applicable, e.g., personally identifiable information or offensive content. Yes
  \end{enumerate}

  \item If you used crowdsourcing or conducted research with human subjects, check if you include:
  \begin{enumerate}
    \item The full text of instructions given to participants and screenshots. No
    \item Descriptions of potential participant risks, with links to Institutional Review Board (IRB) approvals if applicable. No
    \item The estimated hourly wage paid to participants and the total amount spent on participant compensation. No
  \end{enumerate}

\end{enumerate}



\clearpage
\appendix
\thispagestyle{empty}


\end{document}
